% Standalone problem document, generated from the adjacent README.md.
% Compile with XeLaTeX. Mathematical content is maintained in Markdown.
\documentclass[11pt,a4paper]{article}
\usepackage[fontsize=10.5pt]{scrextend}
\usepackage[margin=22mm,headheight=15pt,headsep=9mm,footskip=11mm]{geometry}
\usepackage{fontspec}
\setmainfont{texgyrepagella}[Extension=.otf,UprightFont=*-regular,BoldFont=*-bold,ItalicFont=*-italic,BoldItalicFont=*-bolditalic]
\setsansfont{texgyreheros}[Extension=.otf,UprightFont=*-regular,BoldFont=*-bold,ItalicFont=*-italic,BoldItalicFont=*-bolditalic]
\setmonofont{texgyrecursor}[Extension=.otf,UprightFont=*-regular,BoldFont=*-bold,ItalicFont=*-italic,BoldItalicFont=*-bolditalic,Scale=MatchLowercase]
\usepackage{amsmath,amssymb}
\usepackage{unicode-math}
\setmathfont{texgyrepagella-math.otf}
\usepackage{xcolor,microtype,enumitem,needspace,fancyhdr}
\usepackage{bookmark}
\definecolor{ink}{HTML}{17344B}
\definecolor{accent}{HTML}{197D80}
\definecolor{muted}{HTML}{596773}
\hypersetup{colorlinks=true,linkcolor=accent,urlcolor=accent,pdftitle={IE-18 - The
exact four-step amplification of restarted Anderson
acceleration},pdfauthor={Open Problems in Numerical Linear Algebra}}
\urlstyle{same}
\setlength{\parindent}{0pt}
\setlength{\parskip}{5pt plus 1pt minus 1pt}
\linespread{1.04}
\setlength{\emergencystretch}{3em}
\widowpenalty=10000
\clubpenalty=10000
\displaywidowpenalty=10000
\allowdisplaybreaks[1]
\setlist{nosep,leftmargin=1.5em}
\providecommand{\tightlist}{\setlength{\itemsep}{0pt}\setlength{\parskip}{0pt}}
\providecommand{\pandocbounded}[1]{#1}
\pagestyle{fancy}
\fancyhf{}
\fancyhead[L]{\sffamily\footnotesize\color{muted}OPEN PROBLEMS IN NUMERICAL LINEAR ALGEBRA}
\fancyhead[R]{\sffamily\bfseries\footnotesize\color{accent}IE-18}
\fancyfoot[L]{\parbox[b]{0.9\textwidth}{\sffamily\fontsize{8}{10}\selectfont\color{muted}Editorial ratings; a bounded search cannot certify that a problem remains unsolved.\\Literature check: 2026-09-08}}
\fancyfoot[R]{\sffamily\footnotesize\color{muted}\thepage}
\renewcommand{\headrulewidth}{0.3pt}
\renewcommand{\headrule}{\hbox to\headwidth{\color{accent}\leaders\hrule height \headrulewidth\hfill}}
\makeatletter
\renewcommand{\section}{\Needspace{5\baselineskip}\@startsection{section}{1}{0pt}{12pt plus 2pt minus 2pt}{4pt}{\sffamily\bfseries\large\color{ink}}}
\renewcommand{\subsection}{\Needspace{4\baselineskip}\@startsection{subsection}{2}{0pt}{9pt plus 1pt minus 1pt}{3pt}{\sffamily\bfseries\normalsize\color{ink}}}
\makeatother
\setcounter{secnumdepth}{0}
\begin{document}
{\sffamily\small\color{muted}Linear systems and elimination\par}
\vspace{3pt}
{\raggedright\hyphenpenalty=10000\exhyphenpenalty=10000\sffamily\bfseries\fontsize{21}{24}\selectfont\color{ink}The
exact four-step amplification of restarted Anderson acceleration\par}
\vspace{7pt}
{\sffamily\small\textbf{Difficulty:} challenging\quad\textbf{Importance:} interesting
to the community\par}
\vspace{4pt}
{\color{accent}\hrule height 0.8pt}
\vspace{6pt}
\textbf{Status:} explicit conjecture in the 2025 journal version

Let \(n\ge2\) and let \(M\in\mathbb R^{n\times n}\) be nonzero and
symmetric, with \(1\notin\sigma(M)\). Put \(A=I-M\). Define the
positively homogeneous map \[
R(v)=M\left(v-\frac{v^TAv}{\|Av\|_2^2}Av\right)\quad(v\ne0),\qquad R(0)=0.
\] Two steps of restarted Anderson acceleration with memory one applied
to \(x=Mx+b\) propagate the residual by \(R\). If \(m_1,\ldots,m_n\) are
the eigenvalues of \(M\), is \[
\max_{v\ne0}\frac{\|R(R(v))\|_2}{\|v\|_2}
=
\max_{i\ne j}
\left(\frac{m_im_j(m_j-m_i)}{|m_i(m_i-1)|+|m_j(m_j-1)|}\right)^2?
\] A term with zero denominator is defined as zero: under the
assumptions this can occur only when \(m_i=m_j=0\).

The conjecture says that a largest four-step residual amplification is
attained using only two orthogonal eigenvectors, irrespective of the
dimension. It supplies an exact worst-case factor for this restart
scheme, relevant to accelerated stationary solvers and multigrid.

\subsection{References}\label{references}

O. A. Krzysik, H. De Sterck, and A. Smith, \emph{Asymptotic convergence
of restarted Anderson acceleration for certain normal linear systems},
SISC 47(2025), Conjecture 10, (9),(16),(24),(27)
(\href{https://doi.org/10.1137/24M1672262}{journal};
\href{https://arxiv.org/html/2312.04776v4}{arXiv v4}). Conjecture 10
states the maximum through two-eigenvector nonlinear eigenvalues;
equation(24) evaluates their scalar maximum. The map formulation above
also covers exact termination without an undefined \(\alpha(0)\).

\subsection{Status check --- 2026-09-08}\label{status-check-2026-09-08}

The arXiv history lists v4,12 May 2025, as latest. Its Conjecture 10 and
journal abstract retain the conditional result. Searches for the paper
title, ``restarted Anderson Conjecture 10'', and later
proof/counterexample results found no resolution. Windowed AA, GMRES(1),
and restarted CG use different residual maps. The September 2026
Forsythe paper resolves the CG restart question and is not a solution
claim for the map above.
\end{document}
